\documentclass[a0paper,
landscape,
%portrait,
fontscale=0.332,
%margin=4cm,
%paperwidth=15truecm,
%paperheight=9truecm,
debug,
hyperref={pdfpagelabels=false}]{baposter}

\newcommand\meta{./meta}

\input{\meta/config/commands}

\setdefaultitem{\textcolor{unibablueI}{$\blacktriangleright$}}%
{\textcolor{unibablueII}{$\triangleright$}}%
{\textcolor{unibablueIII}{-}}{\textperiodcentered}

\setdefaultenum{\color{unibablueI}I.}{\color{unibablueII}a)}{i.}{A.}

\newcommand{\colred}[1]{%
\textcolor{nounibaredI}{\textbf{#1}}%
}

\newcommand{\colblue}[1]{%
\textcolor{unibablueI}{#1}%
}

%%%%%%%%%%%%%%%%%%%%%%%%%%%%%%%%%%%%%%%%%
%%%%%%%%%% Glossary entries    %%%%%%%%%%
%%%%%%%%%%%%%%%%%%%%%%%%%%%%%%%%%%%%%%%%%

\newacronym{ieee}{IEEE}{Institute of Electrical and Electronics Engineers}

\newglossaryentry{latex}
{
        name=latex,
        description={Is a mark up language specially suited for scientific documents}
}

\newglossaryentry{symb:Pi}
{
        name=\ensuremath{\pi},
        description={Geometrical value},
        type=symbolslist
}

\lstset{language=TeX}

\begin{document}

\begin{poster}{
  %grid=true,
  background=plain,%shadelr,
  columns=4,
  colspacing=0.2cm,
  linewidth=0.2pt,
  bgColorOne=unibagrayV,
  bgColorTwo=unibaredIII,
  headerheight=0.12\textheight,
  footerheight=0.02\textheight,
  borderColor=unibagrayV,
  headerColorOne=unibayellowII,
  headerColorTwo=unibayellowIV,
  headerFontColor=unibablueI,
  headerborder=none,
  headershape=roundedright,
  headershade=shadetbinverse,%plain,%shadelr,
  headerfont=\large\bf,
  boxColorOne=white,
  %boxpadding=1em,
  %boxheaderheight=2em,
  textborder=roundedleft,
  %eyecatcher=false
}
{}
%%%%%%%%%%%%%%%%%%%%%%%%%%%%%%%%%%%%%%%%%%%%%%%%%%%%%%%%%%%%%%%%%%%%
%%%%                  Metainfo                                  %%%%
%%%%%%%%%%%%%%%%%%%%%%%%%%%%%%%%%%%%%%%%%%%%%%%%%%%%%%%%%%%%%%%%%%%%
% TODO: Ausgliedern
% Title
 {\Huge\textcolor{unibayellowI}{\textbf{\textsc{Title}}}\normalsize}
 % Authors
 {\LARGE\textcolor{unibagrayV}{\textbf{Author}}\normalsize\\[.5em]
 {\Large\textcolor{unibagrayIV}{\textbf{Computer Networks Group, University of Bamberg, Germany}}\normalsize}}
{\includegraphics[width=.07\linewidth]{logow}}
{\textcolor{unibablueI}{\texttt{\url{http://www.uni-bamberg.de/ktr}}}}
{\textcolor{unibablueI}{\textbf{University of Bamberg, Germany}}}
{\textcolor{unibablueI}{\textbf{Email:} \texttt{Email}}}

%%%%%%%%%%%%%%%%%%%%%%%%%%%%%%%%%%%%%%%%%%%%%%%%%%%%%%%%%%%%%%%%%%%%
%%%%                  Content                                   %%%%
%%%%%%%%%%%%%%%%%%%%%%%%%%%%%%%%%%%%%%%%%%%%%%%%%%%%%%%%%%%%%%%%%%%%

\begin{posterbox}[column=0, name=posterboxOne, height=auto]{Integration von Bildern}
Bilder befinden sich im Image-Ordner und lassen sich mit \lstinline|\image{Breite}{Datei im Image-Verzeichnis}{Beschriftung}{Label}| einbinden.
\image{3cm}{logo.png}{Uni-Logo}{img:uni}
Die Referenzierung erfolgt mittels \lstinline|\autoref{Label}|, also z.B. in Abbildung \ref{img:uni}.
\end{posterbox}

\begin{posterbox}[column=1, span=2,height=auto, name=posterboxTwo, boxpadding=1em]{Git Versionierung}
\begin{itemize}
\item \#: \gitAbbrevHash
\item @: \gitAuthorIsoDate
\item \gitReferences
\end{itemize}
\textbf{Setup}\\
Follow this link to \url{https://www.ctan.org/tex-archive/macros/latex/contrib/gitinfo2}\\
\textbf{git hooks}\\
  To fill watermark at buttom, deploy gitinfo2-hook.txt to githooks: (copy and make executable) or use \texttt{make git}
  \begin{itemize}
      \item .git/hooks/post-checkout
      \item .git/hooks/post-commit
      \item .git/hooks/post-merge
  \end{itemize}
  \textbf{Remove watermark}\\
  To disable watermark, remove option \texttt{[mark]} from \textbackslash usepackage[mark]\{gitinfo2\} in \textit{config/commands.tex}.
\end{posterbox}

\begin{posterbox}[column=0,span=2,below=auto,height=bottom, name=posterboxThree]{Grafiken mit TikZ}
Grafiken im TikZ-Framework\footnote{\url{http://www.tn-home.de/TUGDD/Stuff/TikZ_final.pdf}} lassen sich mit dem Befehl \lstinline|\scaletikzimage{Datei im Image Verzeichnis}{Beschriftung}{Label}{Skalierungsfaktor}| einbinden. \scaletikzimage{tikz}{TikZ-Grafik}{img:tikz}{0.9}

Grafiken mit Netzwerksymbolen können einfach mit moeptikz\footnote{\url{https://github.com/moepinet/moeptikz}} erstellt und mit \lstinline|\scaletikzimage{}{}{}| eingebunden werden:
\scaletikzimage{network}{Example Network}{img:topo}{0.9}

UML basierte Grafiken können einfach mit tikz-uml\footnote{\url{https://perso.ensta-paris.fr/~kielbasi/tikzuml/index.php}} erstellt werden und mit \lstinline|\scaletikzimage{}{}{}| eingebunden werden:
\scaletikzimage{sequence}{Example Sequence}{img:seq}{0.9}
\end{posterbox}

\begin{posterbox}[column=2,below=posterboxTwo, height=auto, name=posterboxFour]{Tabellen}
Tabellen, siehe Tabelle \ref{lst:table}, lassen sich mit dem Environment \lstinline|longtable| definieren\footnote{\url{ftp://ftp.dante.de/pub/tex/macros/latex/required/tools/longtable.pdf}}.

\begin{lstlisting}[caption=Tabelle, language=TeX, label=lst:table]
\begin{longtable}[H h t b c]{Spaltendefinitionen} \\
Zelle 1 & Zelle 2 & ... & Zelle n \\
... \\
Zelle x & Zelle y & ... & Zelle z \\
\caption{Tabellenunterschrift}
\label{Label}
\end{longtable}
\end{lstlisting}

\begin{longtable}[H]{|p{0.2\textwidth}|p{0.2\textwidth}|p{0.2\textwidth}|}
\hline
Zelle 1 & Zelle 2 & Zelle n \\
\hline
Zelle x & Zelle y & Zelle z \\
\hline
\caption{Tabelle 1}
\label{tab:tab1}
\end{longtable}
\end{posterbox}

\begin{posterbox}[column=2, height=bottom, below=posterboxFour, name=posterboxFive]{Abkürzungen}
Abkürzungen können mit \lstinline|\newacronym{ieee}{IEEE}{Institute of Electrical and Electronics Engineers}| angegeben werden. Diese werden alphabetisch sortiert in ein Abkürzungsverzeichnis aufgenommen und im Text z.B. mit \lstinline|\gls{ieee}| referenziert, dies führt bei Verwendung zu \gls{ieee}.
\end{posterbox}

\begin{posterbox}[column=3, height=auto, name=posterboxSix]{Code Ausschnitte}
Pseudo-Code Ausschnitte lassen sich mit \lstinline|pseudo{Name des Algorithmus}{Label}{Datei im Code-Verzeichnis}| einbinden.
\pseudo{Mittelwert}{lst:mean}{code}

Code Ausschnitte lassen sich einfach mit listings verwenden z.B. in der Umgebung \lstinline|\begin{lstlisting} ... \end{lstlisting}|, wie in Code Ausschnitt \ref{lst:python} gezeigt.

\begin{lstlisting}[language=Python, caption={Simple Python program}, label=lst:python]
def my_function():
  print("Hello from a function")
\end{lstlisting}
\end{posterbox}

%%%%%%%%%%%%%%%%%%%%%%%%%%%%%%%%%%%%%%%%%%%%%%%%%%%%%%%%%%%%%%%%%%%%
%%%%                  Bibliography                              %%%%
%%%%%%%%%%%%%%%%%%%%%%%%%%%%%%%%%%%%%%%%%%%%%%%%%%%%%%%%%%%%%%%%%%%%

\begin{posterbox}[column=3, below=posterboxSix, height=bottom]{References}
%\renewcommand{\refname}{}
\small
\nocite{*}
\bibliographystyle{abbrv}
\renewcommand{\section}[2]{\vskip 0.05em}
\bibliography{\meta/exampleLiterature/bib}
\end{posterbox}


\end{poster}
\end{document}
